\subsection*{First Reading: 1 Kings 3:5, 7--12}

\subsubsection*{The pericope, and the three things it leaves out}

The scene is Solomon's night at Gibeon, ``the great high place,'' early in a reign that has just been secured by violence. The Lectionary appoints the offer, the request and its approval, and it draws three boundaries that change the reading's shape.

\begin{threecoltable}{Omission}{What stands in the Bible}{Effect on the proclaimed text}
\rowthree{v.~6, inside the appointed span}{Solomon's recital of the great mercy shown to David his father, and of the son given to sit on his throne.}{The dynastic warrant is removed. What the assembly hears is a young man's self-description without the family credential that in the text precedes it.}
\rowthree{vv.~13--14, immediately after}{God adds riches and glory unasked, and makes long life conditional on walking in his ways.}{The reading ends on the granting of the wise heart. The reward that follows the refusal of reward, and the condition attached to long life, are not proclaimed.}
\rowthree{vv.~16--28, later in the chapter}{The judgment between the two women and the living child.}{The famous demonstration is absent. Anything a homily draws from ``the judgment of Solomon'' is drawing on a text that was not read.}
\end{threecoltable}

Verse~3 has already noted that Solomon sacrificed and burnt incense in the high places, and the narrator says so as a qualification, not a compliment. The book will end with Solomon's apostasy and the kingdom's division. Neither fact is proclaimed today, and neither should be suppressed: the man who asks so well is the same man the book will judge, and the request is not a character reference.

\subsubsection*{One phrase, three languages}

The request is usually quoted as ``an understanding heart,'' and that rendering conceals what is asked for.

\begin{threecoltable}{Witness}{1 Kings 3:9, the request}{What is asked}
\rowthree{Hebrew}{\emph{l\=eb sh\=om\=ea{}}, with \emph{sh\=om\=ea{}} a Qal participle of \emph{sh\=ama'}, to hear}{A \emph{hearing} heart: an organ of attention, not of analysis. The verb is rendered ``hear'' or ``hearken'' in the overwhelming majority of its biblical occurrences and ``understand'' in a handful.}
\rowthree{Septuagint (3 Kingdoms 3:9)}{\emph{kardian akouein}}{Literally ``a heart to hear.'' The Greek keeps the Hebrew idiom, then moves at v.~11 to understanding for hearing judgment and at v.~12 to a prudent and wise heart.}
\rowthree{Clementine Vulgate (3 Kings 3:9)}{\incipit{cor docile}}{A \emph{teachable} heart. The Latin shifts the metaphor from listening to docility.}
\rowthree{Douay-Rheims}{``an understanding heart''}{The English follows the Latin's direction and completes the shift toward cognition.}
\end{threecoltable}

What God grants is described differently again: at v.~12 the Latin gives \incipit{cor sapiens et intelligens}, a wise and understanding heart, and at v.~11 the thing asked is summarised as \incipit{sapiéntiam ad discernéndum iudícium}, wisdom to discern judgment. So the request is for hearing, the summary is for discernment, and the gift is for wisdom and understanding. Hebrew, Greek and Latin all perform that movement; the Latin alone begins it already at teachability.

\subsubsection*{Reception, with its limits stated}

Ambrose's treatment in \emph{On the Duties of the Clergy} II.8 is the standard patristic locus, and it needs a precise description. The chapter is not about Gibeon at all: it is about \emph{counsel} as the third thing that commends a man to others, and its thesis is that good advice requires prudence and justice together, since ``prudence cannot exist without justice'' (§43). Ambrose then narrates the judgment between the two women (§§44--45), calls Solomon's penetration of the hidden maternal affection divine (§46), and divides the achievement: it was wisdom ``to distinguish between secret heart-thoughts, to draw the truth from hidden springs,'' and justice that the false claimant should not take another's child (§47). Only in that section's closing clause does the appointed verse appear, and it appears as the judgment's ground: Solomon ``had asked for wisdom, so that a prudent heart might be given him to hear and to judge with justice.'' Ambrose therefore does receive the appointed verse directly, but backwards, from the scene the Lectionary does not appoint.

Augustine's Solomon is a different figure again, and worth setting beside Ambrose's because it cuts the other way. In the \emph{City of God} Solomon is a type who fails to resemble what he typifies: whoever thinks Nathan's promise was fulfilled in him ``greatly errs,'' since one need only look at the house of Solomon ``full of strange women worshipping false gods'' (XVII.8). And the diagnosis in XVII.20 is exact about the appointed scene's aftermath: ``This man, after good beginnings, made a bad end. For indeed prosperity, which wears out the minds of the wise, hurt him more than that wisdom profited him.'' The man who asks so well is not thereby safe, and the book that records the request also records the ruin.

\subsubsection*{Four bounded negatives}

Where a witness would be expected and is absent, that is worth recording, because a homilist who assumes the tradition is full of this text will misjudge what may be claimed for it.

Gregory the Great's \emph{Pastoral Rule} is the obvious candidate: a whole treatise on the government of souls, written by a man who had been a prefect. Reading all four books, Solomon appears sixteen times, every time as the tag attached to a quotation from Proverbs, Ecclesiastes or Sirach, and Gibeon does not appear at all. The nearest approach is Book III's warning that Solomon, granted such wisdom, fell because no discipline of tribulation guarded it -- the grant presupposed, the petition never invoked.

Cassian's second \emph{Conference}, the classic treatise on discretion, likewise never cites the request. Its royal exempla are failures: Saul, who preferred his own offering to obedience, and Ahab, who preferred his own clemency to a divine command. The tradition's foundational text on discernment was built on kings who lacked it rather than on the king who asked for it.

Aquinas is the third. Reading the whole of \emph{Summa theologi\ae} II-II q.~45 on the gift of wisdom, q.~47 on prudence, and q.~83 on prayer, and I-II q.~68 on the gifts of the Spirit, none contains any reference to Solomon or to 3~Kings~3. The absence is most striking at q.~83, a.~6, where Aquinas holds that temporal things may lawfully be desired ``not indeed principally, by placing our end therein, but as helps whereby we are assisted in tending towards beatitude'' -- which is precisely the logic of Solomon's refusal, argued instead from Augustine. The convergence is real and the dependence is not evidenced. This is a bounded result for those four questions, not a claim about the whole Thomistic corpus.

The fourth is modern and belongs in the gallery below rather than here: what Benedict~XVI did with this verse in a national parliament changes the text's register rather than expounding it.

\subsection*{Responsorial Psalm: Ps 119:57, 72, 76--77, 127--128, 129--130}

\subsubsection*{What the selection is made of}

Psalm 119 is the Psalter's longest poem and its most rigidly built: twenty-two stanzas of eight verses, each stanza's verses beginning with the same successive letter of the Hebrew alphabet, and almost every verse containing a synonym for the law. Knowing where the appointed verses sit inside that grid is not antiquarianism; it shows what the Lectionary has done.

\begin{fourcoltable}{Sung verses}{Stanza (letter, ordinal)}{Position within the stanza}{What the verse says, in the Douay-Rheims}
\rowfour{57}{HETH, 8th (vv.~57--64)}{First verse}{``O Lord, my portion, I have said, I would keep thy law.''}
\rowfour{72}{TETH, 9th (vv.~65--72)}{Last verse}{``The law of thy mouth is good to me, above thousands of gold and silver.''}
\rowfour{76--77}{JOD, 10th (vv.~73--80)}{Fourth and fifth}{Mercy for comfort, and tender mercies that one may live, ``for thy law is my meditation.''}
\rowfour{\emph{Response:} 97a}{MEM, 13th (vv.~97--104)}{First verse}{``O how have I loved thy law, O Lord!''}
\rowfour{127--128}{AIN, 16th (vv.~121--128)}{Last two}{Commandments loved ``above gold and the topaz''; ``I have hated all wicked ways.''}
\rowfour{129--130}{PHE, 17th (vv.~129--136)}{First two}{``Thy testimonies are wonderful''; ``the declaration of thy words giveth light: and giveth understanding to little ones.''}
\end{fourcoltable}

Three features of the selection are objective and consequential. It is discontinuous: it draws from six of the twenty-two stanzas and skips the rest. It is edged: it takes the opening verse of one stanza, the closing verse of another, the closing pair of a third and the opening pair of a fourth, so that the sung text repeatedly stands at a seam. And its refrain comes from a stanza none of whose other verses is sung, and which falls chronologically between the second and third of the sung groups. No ancient commentator expounded this sequence as a unit, because the unit is a modern liturgical construction. What the Fathers expounded is the individual verses, and they did so at length.

\subsubsection*{Two Fathers who found the psalm hard}

Both of the great Latin expositors of Psalm 118 leave a record of the difficulty, and both leave a remark about the appointed verses that is worth having.

Augustine expounded every other psalm before this one, and his preface says why: he deferred it ``not so much on account of its well-known length as on account of its depth, knowable to few.'' He adds the reason that difficulty is not the usual kind: with other obscure psalms the obscurity at least shows itself, whereas this one ``presents such a surface that it is believed to need a reader and a hearer, not an expositor.'' Anyone tempted to treat the responsorial psalm as the day's easy element has been warned by the man who wrote thirty-two sermons on it.

Ambrose supplies something more particular. Opening the section that begins at the first of our appointed verses, he records a dispute about where that verse belongs: most codices attach v.~57 to the \emph{end} of the seventh letter, but the Greek psalter corrected against the Hebrew shows it opens the eighth. The Lectionary's first sung verse therefore sits on a seam already contested in the fourth century.

\subsubsection*{Direct exegesis at the appointed verses}

Augustine's exposition happens to cover five of the six appointed groups directly. The paragraph numbers below are those of the standard English edition, which is abridged and does not print his thirty-two sermon divisions; his Latin lemma at v.~130 also differs slightly from the Gallican Psalter the liturgy uses, reading that God's words \emph{make the little ones understand} rather than \emph{give understanding to the little ones}.

\begin{twocoltable}{Verse}{Augustine's comment, at the English edition's paragraph number}
\rowtwo{57, ``my portion''}{The possession is deferred and conditional in a particular way: ``the Lord will be each man's portion then, when he has kept His law'' (§57).}
\rowtwo{72, above gold and silver}{He reads the comparison as one between two loves rather than two valuations, and the sentence is quotable as it stands: ``love loves the law of God more than avarice loves thousands of gold and silver'' (§71).}
\rowtwo{76, mercy for comfort}{He notes that the preceding verse has already said ``in thy truth thou hast humbled me,'' and takes mercy and truth as the pair in which the psalm asks to be consoled (§76).}
\rowtwo{127--128, above gold and topaz}{Grace has this object, ``that the commandments, which could not be fulfilled by fear, may be fulfilled by love''; therefore they are above gold and topaz (§126). The hatred of the unrighteous way is then the navigator's avoidance of the rock on which precious freight is lost (§127).}
\rowtwo{129--130, wonderful, and light to little ones}{Wonder is the reason for search, not a substitute for it: because the testimonies are wonderful, the soul searched them, growing ``more curious from the difficulty'' (§128). And on the little ones: ``What is the little one save the humble and weak?'' -- the law was given ``that it might make you a little one instead of great'' (§129).}
\end{twocoltable}

Two of these bear directly on the day. The comment on v.~72 supplies the psalm's own answer to the first reading: the ranking of instruction over bullion is not asceticism about money but a competition between loves, which is exactly what the Gospel's two buyers demonstrate in action. And the comment on v.~130 lands on the same Latin word as the Gospel acclamation. The psalm says that God's words give understanding to \emph{parvuli}; the acclamation says the Father has revealed these things to \emph{parvuli}. Augustine's gloss -- the humble and weak, made little rather than great -- is the bridge, and it was written about the psalm, not about the pairing.

Ambrose reaches the same two verses by a different route, and his register is social rather than epistemological. On v.~57 he asks how rare on earth is the man who can say that the Lord is his portion, and answers with a portrait: one whom lust does not inflame, avarice does not goad, ambition does not lay low, envy does not waste -- a man ``born for God, not for himself.'' He then works through the Levite, who has no earthly inheritance because God is his portion, to Peter at the Beautiful Gate, whose ``silver and gold have I none'' becomes the psalm verse's proof: this is my portion, and my portion is Christ. On v.~72 he is blunter still. Not everyone says that the law is better than thousands of gold and silver, ``indeed it is a rare man who says it''; Peter said it and proved it by the effect. And Ambrose then names who cannot: not the miser brooding on buried gold, and not the moneyed man who ransacks his daily profits, piles up wealth, lays snares for legacies, and keeps tireless vigil at a sick man's bedside. The legacy-hunter of Roman satire has walked into a psalm commentary, and he is there because the verse is about a ranking of loves that a whole way of life contradicts.

The psalm's own horizon must be kept. These are the words of an Israelite whose delight is in \emph{torah}, and Christian reception does not annul that literal sense. The Old Testament is proclaimed in this liturgy as Scripture in its own right, not as a cipher awaiting decryption, and the two Fathers quoted above both treat it that way even where their conclusions are Christological.
